\section{Implications and Limits}\label{sec:agenda}
\subsection{A deployment evaluation checklist}
Before claiming that a defense protects a deployment, report: (1) the hidden target and protection relation; (2) the final receiver's actual fields and linkage permissions; (3) where protection ends and whether cover reaches action dispatch; (4) protected-data access, fitting procedure and a matched defended-test contrast where adaptation is claimed; (5) the workload and the generator shifts actually tested; and (6) inference, objective/policy utility and applicable cost bound to the same condition. Report missing entries as scoped evidence states. The checklist asks for evidence needed by the specific claim, not every measurement in every paper.

The three cases make this actionable. A content-aware downstream renderer calls for paired final-hop captures; surviving actions call for residual-feature prediction; executed dummy actions call for objective state/goal tests. Prefix-cache isolation needs workload reuse costs, while a semantic matcher needs a separate mechanism check. Task, topology, browser configuration and hardware shifts need their respective bridge tests. These follow from the source comparisons rather than from a global defense ranking.

\subsection{Threats to the synthesis}
\paragraph{Selection and source binding}
The acquired frames and coded collection do not exhaust IEEE/ACM or preprint literature. Remaining source-review and high-risk exclusion work may reveal a closer systematization, a satisfying deployment tuple or a counterexample. Four empirical-source/lineage questions remain pending; B069 has an abstract-supported background disposition with later pages inaccessible. Four quantitative defense result/run bindings and AHB's historical detector-fit binding remain unverified; the corresponding author statements are source-established. Their uncertain quantities are not premises of the three cases. Removing NetEcho withdraws the sole direct re-emission instance; removing AHB withdraws the measured mobile state-change instance. Total corpus size does not increase those replication counts.

\paragraph{Interpretive reliability}
AI-assisted source checking does not establish independent researcher adjudication of eligibility, result splitting or transfer. A quotation, source hash or verified code path establishes its corresponding fact, not a missing experiment. The independent-review packets make these decisions inspectable, but unreturned forms supply no agreement statistic. The negative workflow comparison establishes no review advantage over the active checklist; its constructed task also does not estimate human review performance.

\paragraph{Empirical reach}
The fixed-archive example supports a representation increment under the stated ordering assumption. Matched-task evidence remains inconclusive; ordinary-site capture, semantic-family independence and live utility/cost remain unmeasured. Failed pilot/retail outcomes are retained. None of these outcomes is replaced by a more favorable metric, and the deployment cases depend on original source results rather than this archive experiment.

\paragraph{Disclosure and materials}
AI assisted screening, extraction and writing; authors remain responsible for the judgments. The companion package contains source-linked tables, saved predictions, protocols and verification code. It excludes private logs, conversations and redistributed third-party datasets. Submission-side access, rights and institutional disclosures require author confirmation; no public repository or external replication is asserted.
